\documentclass[12pt]{article}
\usepackage[T2A]{fontenc}
\usepackage[utf8]{inputenc}
\usepackage[russian]{babel}
\usepackage{latexsym,amssymb,amsthm}
\usepackage{amsmath}
\usepackage{wrapfig}
\RequirePackage[pdftex]{graphicx}
%%\usepackage{array}

\usepackage{epsf,epic,eepic}
\input ris.tex
\input cells.sty
% \addrisoption{\risleftfalse}

\font\Chess=chess10 scaled 600

\DeclareRobustCommand{\No}{\ifmmode{\nfss@text{\textnumero}}\else\textnumero\fi}

\oddsidemargin=0pt
\textwidth=159mm
\headheight=0pt
\headsep=0pt
\topmargin=0pt
\textheight=240mm

%%
%% Три точки для обозначения делимости
%%
\catcode`\@=11
\def\kratno{\mathrel{\smash{\lower.5ex\hbox{$\,\vdots\,$}}}}
% Это прямые три точки для обозначения делимости
\def\nekratno{\mathrel{\mathpalette\c@ncel\kratno}}
\def\c@ncel#1#2{\m@th\ooalign{$\hfil#1\mkern1mu/\hfil$\crcr$#1#2$}}
% Это перечеркнутые три точки для обозначения неделимости
\def\divis{\smash{\lower.1ex\hbox{\,\hbox{.}\kern-.22em\lower-.6ex\hbox{.}
\kern-.52em\lower-1.2ex\hbox{.}\,}} }
% А это наклонные...

\hyphenation{чис-ло чис-ла чис-лу чис-лом чис-ле чис-лам чис-ла-ми чис-лах
че-ты-рех-уголь-ник че-ты-рех-уголь-ни-ка пря-мо-уголь-ник тре-уголь-ни-ка
тре-уголь-ни-ков вне-впи-сан-ной вне-впи-сан-ных}

\sloppy
\pagestyle{empty}

\def\q#1.{{\bf #1. }}
\def\dotline{\smallskip\hbox to \hsize{\dotfill}\medskip}

\begin{document}

\centerline{\sc САНКТ-ПЕТЕРБУРГСКАЯ}
\centerline{\sc ОЛИМПИАДА ШКОЛЬНИКОВ 2023 года ПО МАТЕМАТИКЕ}
\centerline{\sc II тур. \ 6 класс.}
\medskip
\hrule
\bigskip

\q1. Несколько пиратов поделили сокровища: каждому досталось пять драгоценных камней общей стоимостью 100\,000 пиастров. Оказалось, что, какого пирата ни возьми, у него есть три камня, которые в сумме стоят меньше, чем вместе стоят два каких-то камня, принадлежащие двум другим пиратам. Докажите, что у~кого-то из пиратов имеется камень, который стоит больше 25\,000 пиастров.

\q2. По кругу стоят 100 детей в синих шапках и 200 в~красных. Известно, что мальчики в синих шапках и девочки в красных шапках говорят правду, а остальные дети лгут. Каждый мальчик сказал: \emph{Все мои соседи в красных шапках}. Каждая девочка сказала: \emph{Все мои соседи в синих шапках}. Сколько мальчиков может быть среди этих детей?

\q3. Назовем \emph{интересным} натуральное число $n$, обладающее следующим свойством: если взять любое натуральное число $a$, делящееся на $n$, и между какими угодно двумя его цифрами вставить три нуля, то получится число тоже делящееся на $n$. Например, число 2 является интересным. Найдите наибольшее интересное число, не делящееся на 10.

\q4. По кругу стоят 100 стаканов, в каждом~--- вода или берёзовый сок. Известно, что нет трех стоящих подряд стаканов с одной и той же жидкостью. За одну пробу можно отпить жидкость из любого стакана и понять, что в нем налито. Можно ли за 75 проб определить содержимое всех стаканов?

\dotline

\centerline{Олимпиада 2023 года. II тур. 6 класс. Выводная аудитория.}
\smallskip

\q5. Незнайка рассказывает, что однажды он расставил в клетках таблицы $20\times 20$ числа от 1 до 400, а затем стал по одному их стирать. За один ход он выбирал одну линию (строку или столбец), в которой осталось хотя бы одно число, и стирал наименьшее число в этой линии. После нескольких ходов в таблице остались числа от 1 до~$k$ и~только они. При каком наименьшем $k$ описываемое Незнайкой событие {\sc не} могло произойти?

\q6. На экране компьютера написано натуральное число~$A$, большее 1\,000\,000. Каждую секунду компьютер вычитает из $A$ число, получающееся из $A$ отбрасыванием двух последних цифр, и заменяет на экране число $A$ на полученную разность. Когда $A$ становится меньше 1000, компьютер останавливается. 
Могло ли так случиться, что во время работы компьютера число $A$ ни в какой момент не оканчивалось двумя нулями?


\clearpage

\centerline{\sc САНКТ-ПЕТЕРБУРГСКАЯ}
\centerline{\sc ОЛИМПИАДА ШКОЛЬНИКОВ 2023 года ПО МАТЕМАТИКЕ}
\centerline{\sc II тур. \ 7 класс.}
\medskip
\hrule
\medskip

\q1. Клетки доски $10\times 10$ покрашены в два цвета. Блоха умеет перепрыгивать с любой клетки или в соседнюю по горизонтали клетку того же цвета, или в соседнюю по вертикали клетку другого цвета. Известно, что такими прыжками блоха может допрыгать с любой клетки доски на любую другую клетку. Докажите, что клеток обоих цветов на доске поровну.

\q2. Школьники писали олимпиаду, за которую можно было получить целое число баллов от 0 до 70. Оказалось, что нулевых работ в 5 раз больше, чем работ с 70 баллами. Оргкомитет выкинул все нулевые работы, отчего средний балл (отношение суммы всех баллов к количеству всех работ) увеличился. Потом жюри заметило, что одна седьмая часть работ с 70 баллами списана, и~оценку в них изменили с 70 на 0. В результате средний балл вернулся к~первоначальному значению. Докажите, что не менее 400 работ имело оценку 1~балл.

\q3. В выпуклом пятиугольнике $ABCDE$ длины всех сторон равны 1. Внутри пятиугольника нашлась такая точка $F$, что $AF=DF$ и $BF=CF=EF$. Найдите расстояние от точки $E$ до прямой $CF$.

\q4. На доске написано число $100!^{100}$. Двое делают ходы по очереди. За ход можно стереть с доски любое число и выписать вместо него два натуральных числа, в произведении дающие стёртое. При этом нельзя выписывать на доску квадраты натуральных чисел. Проигрывает игрок, который не может сделать ход. Кто выиграет при правильной игре --- начинающий или его соперник?

\dotline

\centerline{Олимпиада 2023 года. II тур. 7 класс. Выводная аудитория.}
\smallskip

\q5. В Тихом Омуте может водиться два вида рыб: мешкороты и
живоглоты. Когда одна рыба съедает другую, у нее вырастает
одна колючка. Рыба может съесть рыбу того же вида, если у
этих двух рыб в сумме четное число колючек; рыба может
съесть рыбу другого вида, если у них в сумме нечётное число колючек.

Ученые запустили в пустой Омут новорожденных рыб без единой колючки: 123 мешкоротов и 321 живоглотов. Может ли
через некоторое время в Тихом Омуте остаться всего одна рыбина?

\q6. В равнобедренном треугольнике проведена биссектриса $AK$. На основании $AC$ отмечена точка $M$, на продолжении отрезка $AK$ за точку $K$ --- точка~$L$, а на продолжении стороны $AB$ за точку $A$ --- точка $N$, причем $AC=CL$, $CK=CM$ и $AN=AM$. Докажите, что $KL = MN$.

\q7. Будем называть расстановку натуральных чисел по кругу
\textit{разделённой}, если в каждой паре соседних чисел одно из чисел делится на другое. Дима записал по кругу $N$ натуральных чисел, образующих разделенную расстановку, в которой любые два соседних числа различны. 
Оля записала между каждыми двумя соседними Димиными числами модуль их разности, а исходные числа стёрла. Полученная Олей расстановка тоже оказалась разделённой. 
Какое наибольшее количество различных чисел могло быть в Диминой расстановке?


\clearpage

\centerline{\sc САНКТ-ПЕТЕРБУРГСКАЯ}
\centerline{\sc ОЛИМПИАДА ШКОЛЬНИКОВ 2023 года ПО МАТЕМАТИКЕ}
\centerline{\sc II тур. \ 8 класс.}
\medskip
\hrule
\bigskip

\q1. На основании $AD$ трапеции $ABCD$ нашлась такая точка $E$, что $AE = EC$ и $BE = ED$. Докажите, что точка пересечения диагоналей трапеции лежит на биссектрисе угла $BEC$.

\q2. Найдите все такие натуральные $n$, что десятичная запись числа $n^2$ начинается с числа $n$. 

\q3. По кругу стоят 300 стаканов, некоторые из них --- с~водой, а~остальные с березовым соком, причём нет трех стаканов подряд с~одинаковой жидкостью. За ход можно отпить из одного стакана. Можно ли за 200 проб гарантированно определить содержимое всех стаканов? 

\q4. Можно ли так расположить на плоскости четыре квадрата попарно разных размеров, чтобы любые два квадрата имели общую вершину, а никакие три общих вершин не имели?

\dotline

\centerline{Олимпиада 2023 года. II тур. 8 класс. Выводная аудитория.}
\smallskip

\q5. У Саши есть 80 кусков пластилина. За один ход он может склеить три любых куска в один, после чего разделить этот кусок на три одинаковых по весу куска. Докажите, что Саша может добиться того, чтобы через несколько ходов среди имеющихся у него кусков пластилина не нашлось бы трех кусков попарно разного веса.

\q6. Дано натуральное число $m$. Сумма нескольких натуральных чисел равна $m^2$. Докажите, что количество чисел, которые являются делителем хотя бы одного из них, меньше чем $2m$.

\q7. Существует ли выпуклый многоугольник, который можно разрезать непересекающимися (во внутренних точках) диагоналями на треугольники равной площади хотя бы тремя разными способами?

\end{document}
